% conclusion.tex - Conclusion

\section{Conclusion}
\label{sec:conclusion}

\lambdaCap{} is a small calculus aimed at three gaps in prior capability
work: the lack of a compact static account of capability use, the lack of a
clean formal model for fine-grained capability propagation through
higher-order code, and the lack of a corresponding foundation for simulation
and testing with explicit capability values. %

The model addresses these gaps with contextual capabilities: capabilities
that are tracked statically, yet may still be supplied by the calling context
rather than captured unconditionally. %
This makes it possible to reason about
two important patterns at once:
\begin{itemize}[nosep]
\item packaging attenuated authority into narrow interfaces, and
\item preserving caller control over invocation-time capabilities such as
  output channels or tool bindings. %
\end{itemize}

The soundness theorem shows that capability annotations are sound with
respect to evaluation. %
If the type system says a term uses no capabilities,
then evaluation succeeds with no capabilities provided. %
More generally, the
model supports rigorous reasoning about authority restriction, capability
containment, and higher-order confinement. %

The result is a verified semantic core for statically checked contextual
capabilities, especially relevant when executing untrusted or LLM-generated
code. %
Its main purpose is to guide the design of new secure programming
languages that make capability use explicit, statically checkable, and
controlled by program context. %

Resource bounds and covert channels are not addressed by this core calculus;
they must be handled by other means, such as runtime resource enforcement and
side-channel controls. %
